% ch10_a §10.1–10.3 (书页 329–335 / p343–p349)
\chapter{磁性}

\epigraphbook{万物皆始于秩序，故亦将终于秩序，并将依循秩序的制定者、\\天上之城的神妙数学，再度开始。}{SIR THOMAS BROWNE}

\section{轨道磁化率}

在本章中，我们所关注的是固体的磁化率本身，而不涉及磁场对固体其他性质（例如电导率）的影响。这是一个很大的课题，我们无法细致处理。例如，我们将略去大多数磁共振（magnetic resonance）现象。我们也不讨论技术上十分重要的、可永久磁化的材料的那些宏观性质的理论，例如磁滞（hysteresis）、矫顽力（coercivity）、剩磁（remanent magnetism）等等；这些性质本质上取决于样品被自发磁化而分成若干磁畴（domains），各磁畴的磁矩取向彼此不同。我们甚至不打算考虑划分磁化不同的磁畴的 Bloch 壁（Bloch wall）的优美理论——这种壁在施加磁场时还可以被驱使移动。

我们首先考虑\emph{非磁性固体}，在其中完全不必怀疑原子或离子携带有局域磁矩。我们也忽略电子自旋，或者假定所有自旋都紧紧地配对。即便如此，电子的轨道运动仍会对磁化率有贡献。

最简单的情形是每个原子或离子都由电子的闭合壳层组成，且激发到更高状态的能量很大。众所周知，这就导致抗磁性（diamagnetism）；磁化率为负，近似地由下式给出：
\begin{equation*}
\chi=-\frac{Ze^{2}N}{6mc^{2}}\,\overline{r^{2}},
\tag{10.1}
\end{equation*}
其中 $Z$ 是原子中电子的总数，$N$ 是单位体积中的原子数，$\overline{r^{2}}$ 是每个原子周围电子电荷云的方均半径。这正是我们处理稀有气体固体、或者离子晶体中离子的贡献时应当使用的那类公式。

上述结果本质上来自一阶微扰能量
\begin{equation*}
\left\langle 0\,\left|\frac{e^{2}}{2mc^{2}}A^{2}\right|\,0\right\rangle,
\tag{10.2}
\end{equation*}
其中 $\mathbf{A}$ 是磁场的矢量势，$|0\rangle$ 是每个原子的基态。当我们按照磁场中哈密顿量的标准表达式，把动能写成
\begin{equation*}
T=\frac{1}{2m}\left(\frac{\hbar}{i}\nabla-\frac{e}{c}\mathbf{A}\right)^{2}
\tag{10.3}
\end{equation*}
的形式时（如式 (9.51) 那样），上式便自动出现。

但是，如果原子存在激发态，比如说 $|n\rangle$，其轨道矩不为零，那么 (10.3) 中线性于 $\mathbf{A}$ 的项——也就是给出算符
\begin{equation*}
\frac{e}{2mc}\,\mathbf{H}\cdot\mathbf{L}=\mathbf{H}\cdot\boldsymbol{\mu}_{L},
\tag{10.4}
\end{equation*}
的那一项——就可能把这些激发态混入基态，其中 $\mathbf{L}$ 是电子的轨道角动量算符，而 $\boldsymbol{\mu}_{L}$ 是与之成比例的通常的磁矩算符。在二阶微扰中，这一项对能量的贡献为
\begin{equation*}
\delta\mathcal{E}=\sum_{n}\frac{|\langle 0|\mathbf{H}\cdot\boldsymbol{\mu}_{L}|n\rangle|^{2}}{\mathcal{E}_{0}-\mathcal{E}_{n}}
\tag{10.5}
\end{equation*}
它是正的，并且是 $H$ 的二次式。这就是原子、离子或分子的恒定的 van Vleck 顺磁性（van Vleck paramagnetism）的来源。

实践表明，把 (10.2) 和 (10.5) 应用于实际固体极其复杂；这就是我们只在原则上引用这些公式的原因。但我们从 (10.5) 中注意到，当激发能量变小时——例如在填满的能带之上有小的带隙时——这一贡献应当增大。当这个带隙趋于零，也就是当我们拥有自由载流子时，会发生什么呢？

这时理论又变得非常复杂，但对自由电子却存在精确解。在 §9.6 中我们导出了磁场中自由电子的能级；只要贯彻那套导致 de Haas--van Alphen 效应的 §9.7 的形式框架，算出整个电子气能量的平均变化并不困难。我们不寻求自由能中的振荡项，而是集中于那些作为 $H$ 的单调函数的项。例如，在 (9.80) 中我们略去了
\begin{equation*}
\frac{1}{4\pi^{2}s^{2}kT}\left[f^{0}(\mathcal{E})\frac{\partial\mathcal{E}}{\partial x}\right]_{x=0}
=\frac{1}{4\pi^{2}s^{2}kT}\,\frac{e\hbar H}{mc}\left[f^{0}(\mathcal{E})\right]_{x=0}
\tag{10.6}
\end{equation*}
（其中利用了 (9.82) 和 (9.58)）。

但这是 (9.79) 中的一项，而 (9.79) 中含有对所有 $s$ 值的求和以及对 Fermi 球的所有切片 $k_{z}$ 的积分。于是，对自由能的总贡献（除去 (9.79) 中与 $H$ 无关的前两项）便是
\begin{align*}
\Delta F_{L} &=-2kT\sum_{s=1}^{\infty}(-1)^{s}\int_{-\infty}^{\infty}\frac{eH}{2\pi^{2}\hbar c}\,\frac{1}{4\pi^{2}s^{2}kT}\,\frac{e\hbar H}{mc}\left[f^{0}(\mathcal{E})\right]_{x=0}dk_{z}\\
&=-\frac{1}{4\pi^{4}}\left(\frac{eH}{c}\right)^{2}\frac{1}{m}\int_{-\infty}^{\infty}\left[f^{0}(\mathcal{E})\right]_{x=0}dk_{z}\sum_{s=1}^{\infty}\frac{(-1)^{s}}{s^{2}}.
\tag{10.7}
\end{align*}
对 $k_{z}$ 的积分本质上量度的是 k 空间中沿 Fermi 球轴线的填充区域的长度，而对 $s$ 的求和给出 $-\frac{1}{12}\pi^{2}$。把它表示成磁化率，我们便得到 \emph{Landau 抗磁性}（Landau diamagnetism）
\begin{equation*}
\chi_{L}=-\frac{e^{2}k_{F}}{12\pi^{2}mc^{2}}.
\tag{10.8}
\end{equation*}

这个结果为负，是因为能级的聚束倾向于增大系统的总能量。从图 171 可以看到这一点。只有当 费米能级恰好位于两个量子化能级的正中间时，电子才能够“凝聚”到各自相应的管道上而不改变其平均能量。当一个量子化能级逼近 费米能级时，它会把电子一同带上去；这就是 $\Delta F_{L}$ 为正的原因。磁化率与温度无关，因为这是一种平均效应，并不要求能级间隔 $\hbar\omega_{H}$ 大于 $kT$。

按照 §9.7 论证的思路，对任意 费米面这一更一般的情形来计算 (10.7) 并不困难。但这样得到的结果并不正确。周期晶格中电子的 Landau 抗磁性需要更复杂的分析，我们在此不拟给出。此外，在 van Vleck 贡献与 Landau 贡献之间还存在交叉项。

\section{自旋顺磁性}

共享同一轨道态的、自旋相反的电子之间的简并，会被磁场解除。在金属中，这引起电子在两种自旋取向之间重新分布，从而产生一个磁矩。这个效应的计算是初等的。假设我们区分 + 自旋与 $-$ 自旋的电子（相对于 $\mathbf{H}$ 的方向）。那么它们的能量将为
\begin{equation*}
\left.\begin{aligned}
\mathcal{E}_{\mathbf{k}+} &= \mathcal{E}(\mathbf{k})-\mu_{0}H,\\
\mathcal{E}_{\mathbf{k}-} &= \mathcal{E}(\mathbf{k})+\mu_{0}H,
\end{aligned}\right\}
\tag{10.9}
\end{equation*}
其中 $\mathcal{E}(\mathbf{k})$ 是不存在磁场时的能量，而 $\mu_{0}$ 是电子的磁矩。

每个状态中的电子数将由两个不同的 Fermi--Dirac 分布给出，如同式 (4.8) 那样，但具有相同的化学势 $\zeta$。于是，正如我们在图 175 中看到的，+ 自旋态容纳的电子比先前要多——总数为
\begin{equation*}
n_{+}=\int_{0}^{\infty}\frac{1}{2}\,\mathcal{N}(\mathcal{E})\,f^{0}(\mathcal{E}_{\mathbf{k}+})\,d\mathcal{E},
\tag{10.10}
\end{equation*}
这是因为在原来能量 $\mathcal{E}$ 处的态密度，是在 + 自旋与 $-$ 自旋的电子之间平分的。

%FIG{175}: {磁场中“上”自旋与“下”自旋的 费米分布。}

对于负自旋电子的密度有相应的公式。两者之差给出一个磁矩
\begin{align*}
M &=\mu_{0}(n_{+}-n_{-})\\
&=\mu_{0}\int_{0}^{\infty}\tfrac{1}{2}\{f_{0}(\mathcal{E}-\mu_{0}H)-f_{0}(\mathcal{E}+\mu_{0}H)\}\,\mathcal{N}(\mathcal{E})\,d\mathcal{E}\\
&\approx\mu_{0}^{2}H\int_{0}^{\infty}\left(-\frac{\partial f^{0}}{\partial\mathcal{E}}\right)\mathcal{N}(\mathcal{E})\,d\mathcal{E}\\
&=\mu_{0}^{2}H\,\mathcal{N}(\mathcal{E}_{F})
\tag{10.11}
\end{align*}
由式 (4.18)。

这个正的磁矩导致 Pauli 顺磁性（Pauli paramagnetism）。它不依赖于温度，并且直接量度 费米能级处的态密度。因此，原则上它应当与电子比热（§4.7）直接相关联；不过，当人们对多电子相互作用作出修正之后（§5.8），这种严格的等价性便不复存在。在实践中做这类比较的困难在于，抗磁性修正 (10.8) 并非可以忽略，而人们对它又知道得不够准确。不过在有些情形下，可以通过观察传导电子的电子自旋共振，直接测量自旋磁化率。

对自由电子气，容易证明
\begin{align*}
\chi_{P} &=\mu_{0}^{2}\,\mathcal{N}(\mathcal{E}_{F})\\
&=\frac{e^{2}k_{F}}{4\pi^{2}mc^{2}},
\tag{10.12}
\end{align*}
它恰好是 Landau 项 (10.8) 数值的三倍。但是，晶格对这两类磁化率的影响却截然不同。例如，假设 费米面是球形的，但“有效质量”变为 $m^{*}$。容易证明，$\chi_{P}$ 依赖于态密度，因而正比于 $m^{*}$；而 $\chi_{L}$ 由电子的轨道运动决定，不经过 Bohr 磁子的介入，所以将\emph{反比}于 $m^{*}$。

与自旋磁化率相联系的另一个现象是 Knight 移动（Knight shift）。在磁场 $\mathbf{H}$ 中，电子自旋被极化，每个电子具有平均磁矩
\begin{equation*}
\langle\boldsymbol{\mu}_{e}\rangle=\frac{1}{n}\,\chi_{P}\,\mathbf{H}.
\tag{10.13}
\end{equation*}
位于 $\mathbf{r}$ 处的一个电子，将与位于 $\mathbf{R}$ 处的原子核的磁矩 $\boldsymbol{\mu}_{I}$ 通过“接触相互作用”（contact interaction）
\begin{equation*}
\mathcal{H}_{\text{int.}}=\tfrac{8}{3}\pi\,(\boldsymbol{\mu}_{I}\cdot\boldsymbol{\mu}_{e})\,\delta(\mathbf{r}-\mathbf{R}),
\tag{10.14}
\end{equation*}
发生作用，这个效应在核共振理论中是众所周知的。把 (10.13) 代入 (10.14)，我们得到一个平均能量
\begin{align*}
\langle\mathcal{H}_{\text{int.}}\rangle &=\tfrac{8}{3}\pi\,\frac{\chi_{P}}{n}\,H\,|\psi(0)|^{2}\,\mu_{I}\\
&=\Delta\mu_{I}\cdot\mathbf{H},
\tag{10.15}
\end{align*}
姑且这样记，其中 $|\psi(0)|^{2}$ 是波函数在核处的模方。这个能量被解释为核的表观磁矩出现一个改变量 $\Delta\mu_{I}$——即它在磁场中共振频率的移动。

显然，Knight 移动应当正比于 $\chi_{P}$。但还有一个因子 $|\psi(0)|^{2}$，它无法由任何独立的实验给出。原则上这只涉及 费米面上的态，因此人们说它“量度了 费米能级处波函数中 $s$ 型成分的多少”。Knight 移动是一个容易测定的好物理量，但它的结果却不那么容易解释。

\section{居里--外斯定律与铁磁性}

现在我们假设每个原子的行为像一个小磁体，磁矩为 $\boldsymbol{\mu}$，例如可以来自过渡金属或稀土元素的离子中 $d$ 电子或 $f$ 电子的未满壳层。在磁场 $\mathbf{H}$ 中，每个小磁体将具有能量 $-\boldsymbol{\mu}\cdot\mathbf{H}$。假设每个原子都与近邻无关；由于它是局域的，它满足 Boltzmann 统计。磁矩为 $\boldsymbol{\mu}$ 的原子所占的比例将正比于
\begin{equation*}
n(\boldsymbol{\mu})=e^{\boldsymbol{\mu}\cdot\mathbf{H}/kT}.
\tag{10.16}
\end{equation*}

如果我们假设每个磁体都可以自由转动，那么平均磁矩必定是
\begin{equation*}
\langle\boldsymbol{\mu}\rangle=\int\boldsymbol{\mu}\,e^{\boldsymbol{\mu}\cdot\mathbf{H}/kT}d\Omega\,\Big/\int e^{\boldsymbol{\mu}\cdot\mathbf{H}/kT}d\Omega,
\tag{10.17}
\end{equation*}
其中 $d\Omega$ 是 $\boldsymbol{\mu}$ 转动时的立体角元。单位体积内 $N$ 个这种原子的集合将具有磁化率
\begin{align*}
\chi &=N\frac{\partial\langle\boldsymbol{\mu}\rangle}{\partial\mathbf{H}}=\int(N/kT)\,\boldsymbol{\mu}\boldsymbol{\mu}\,e^{\boldsymbol{\mu}\cdot\mathbf{H}/kT}d\Omega\,\Big/\int e^{\boldsymbol{\mu}\cdot\mathbf{H}/kT}d\Omega\\
&\approx\frac{N}{kT}\langle\boldsymbol{\mu}\boldsymbol{\mu}\rangle=\frac{1}{3}\,\frac{N\langle\mu^{2}\rangle}{kT}
\tag{10.18}
\end{align*}
此时 $H$ 已小到指数项可以忽略；并且我们已经把并矢张量 $\boldsymbol{\mu}\boldsymbol{\mu}$ 的平均，同矢量 $\boldsymbol{\mu}$ 之长度的平方 $\mu^{2}$ 的平均作了比较。

如今对我们来说，更自然的做法是把每个磁体同一个自旋角动量 $\mathbf{S}$ 联系起来，它沿磁场方向量子化：
\begin{equation*}
\boldsymbol{\mu}\cdot\mathbf{H}=\beta HS_{z},
\tag{10.19}
\end{equation*}
其中 $S_{z}$ 以整数台阶从 $S$ 变到 $-S$，而 $\beta$ 是 Bohr 磁子的两倍。例如，设 $S=\frac{1}{2}$。这时在 (10.17) 中我们遇到的是求和而不是积分：
\begin{align*}
\langle\boldsymbol{\mu}\rangle &=\beta S\left[\frac{e^{\beta SH/kT}-e^{-\beta SH/kT}}{e^{\beta SH/kT}+e^{-\beta SH/kT}}\right]\\
&=\beta S\,\tanh\frac{\beta SH}{kT}.
\tag{10.20}
\end{align*}

在 $H$ 很小的极限下，这个结果便与 (10.18) 相同，只要我们注意到，按照总角动量大小的规则，
\begin{equation*}
\langle\mu^{2}\rangle=S(S+1)\beta^{2}.
\tag{10.21}
\end{equation*}

结果 (10.18) 就是著名的顺磁磁化率的 \emph{Curie 定律}（Curie Law）。某些类型的固体近似地服从它，尤其是磁性原子或离子在晶体中彼此相距很远的那些固体，因为这时独立性假设成立。

然而在许多物质中，近邻自旋之间存在相互作用。为了近似地把这一点考虑进去，可以引入 \emph{Weiss 场}（Weiss field）
\begin{equation*}
\mathbf{H}_{I}=\lambda N\langle\boldsymbol{\mu}\rangle.
\tag{10.22}
\end{equation*}
这个场被认为来自每个原子所处的磁化平均环境——不过，用来量度相互作用强度的参量 $\lambda$，我们让它保持任意。

作用在一个原子上的总场现在是 $\mathbf{H}+\mathbf{H}_{I}$。把它代替 $\mathbf{H}$ 代入 (10.20)，我们近似地得到
\begin{equation*}
N\langle\boldsymbol{\mu}\rangle=\frac{1}{3}\,\frac{N\langle\mu^{2}\rangle}{kT}\left(\mathbf{H}+\lambda N\langle\boldsymbol{\mu}\rangle\right),
\tag{10.23}
\end{equation*}
由此得
\begin{align*}
\chi &\equiv N\left\langle\frac{\partial\boldsymbol{\mu}}{\partial\mathbf{H}}\right\rangle\\
&=\frac{N\langle\mu^{2}\rangle}{3k(T-\theta)},
\tag{10.24}
\end{align*}
其中
\begin{equation*}
\theta=\frac{\lambda N\langle\mu^{2}\rangle}{3k}.
\tag{10.25}
\end{equation*}
这就是所谓的\emph{居里--外斯定律}（Curie--Weiss law），它描述许多固体的磁化率行为。

参量 $\theta$ 显然具有温度的量纲。如果在 (10.24) 中 $T<\theta$ 会怎么样？磁化率似乎会先变为无穷大，然后变为负值。我们必须回到 (10.20)，求解方程
\begin{equation*}
N\langle\boldsymbol{\mu}\rangle=N\beta S\,\tanh\left\{\frac{\beta S}{kT}\left(\lambda N\langle\boldsymbol{\mu}\rangle+H\right)\right\}.
\tag{10.26}
\end{equation*}
当 $H=0$ 时，通常的解是 $\langle\boldsymbol{\mu}\rangle=0$；但当 $T<\theta$ 时，我们发现还有另一个 $\langle\boldsymbol{\mu}\rangle\neq 0$ 的根。这个根很容易用图解法定位；在温度 $T=0$ 时它显然趋于
\begin{equation*}
\mathbf{M}(0)=N\beta S.
\tag{10.27}
\end{equation*}
